
\begingroup
\shorthandoff{<>}% Babel español trata de otro modo ciertos delimitadores de Xy-pic.
\setcounter{proposition}{10}% Conservar el número visible XIII.2.11 de la fuente.
\begin{lemme}
\label{XIII.2.11}
Sean $f\colon X\to S$ un morfismo propio de presentación finita,
$D$ un divisor sobre~$X$ de cruzamientos normales relativamente a~$S$,
$Y=\Supp D$, $U=X-Y$, $i\colon U\to X$ el morfismo canónico,
$\bar{s}_1\to\bar{s}_2$ un morfismo de especialización de puntos
geométricos de~$S$, $y_1$ un punto geométrico de
$Y_{\bar{s}_1}$, $y_2$ un punto geométrico de
$Y_{\bar{s}_2}$,
\ifthenelse{\boolean{orig}}
{de modo que la proyección $z_1$ de~$y_1$ sobre~$X$ sea una
generalización de la proyección $z_2$ de~$y_2$.}
{de modo que la proyección $z_1$ de~$y_1$ sobre~$X$ sea una
generalización de la proyección $z_2$ de~$y_2$; elijamos además un
morfismo de especialización $y_1\to y_2$ compatible con ella y
supongamos que $z_1$ y $z_2$ pertenecen a los mismos componentes
irreducibles de~$D$.}
Sea $I_{y_1}$
un subgrupo de inercia de~$\pi_1^{\tame}(\bar{U}_{\bar{s}_1})$ en
$y_1$. Entonces la imagen de~$I_{y_{1}}$ por el homomorfismo de
especialización
$$
\pi\colon \pi_1^{\tame}(U_{\bar{s}_1})\to
\pi_1^{\tame}(U_{\bar{s}_2})
$$
es
\marginpar{392}
un subgrupo de inercia de~$\pi_1^{\tame}(U_{\overline{s}_2})$
en~$y_2$.
\end{lemme}

En efecto, sean $\overline{X}$ (\resp $\widetilde{X}$) la localización
estricta de~$X$ en $y_2$ (\resp en $y_1$),
$\overline{U}=U \times_X \overline{X}$
(\resp $\widetilde{U}=U\times_X \widetilde{X}$). Hay un morfismo
canónico $\widetilde{U}\to \overline{U}$, y de~
\ifthenelse{\boolean{orig}}
{\Ref{XIII.1.10}}
{\Ref{XIII.2.10}}
resulta el diagrama conmutativo
$$
\xymatrix{ \pi_1^{\tame} \left(\widetilde{U}_{\overline{s}_1} \right)
\ar[r]^{\pi '} \ar[d] & \pi_1^{\tame} \left(
\overline{U}_{\overline{s}_2} \right) \ar[d] \\ \pi_1^{\tame} \left(
U_{\overline{s}_1} \right) \ar[r]^{\pi } & \pi_1^{\tame} \left(
U_{\overline{s}_2} \right) }
$$
donde $\pi'$ es el compuesto del homomorfismo canónico
$\pi_1^{\tame}(\widetilde{U}_ {\overline{s}_1})\to
\pi_1^{\tame}(\overline{U}_{\overline{s}_1})$, seguido del homomorfismo de
especialización. Como
$\pi_1^{\tame}(\widetilde{U}_{\overline{s}_1})$
(\resp
\ifthenelse{\boolean{orig}}
{$\pi_1^{\tame}(\overline{U}_{\overline{s}_1})$}
{$\pi_1^{\tame}(\overline{U}_{\overline{s}_2})$})
es un grupo de inercia de
$\pi_1^{\tame}(U_{\overline{s}_1})$ en~$y_1$
(\resp de~$\pi_1^{\tame}(U_{\overline{s}_2})$ en~$y_2$), basta
demostrar que $\pi'$ es sobreyectivo. Pero esto resulta de la expresión
obtenida en~\Ref{XIII.5.6}.

\endgroup
